\documentclass[11pt]{amsproc}

\usepackage[a4paper,margin=2.6cm]{geometry}
\usepackage{amsmath,amssymb,amsthm}
\usepackage[T1]{fontenc}
\usepackage[utf8]{inputenc}

\newcommand{\id}{\mathrm{id}}

\begin{document}

\section*{Problema}

Sean $n$ y $m$ enteros positivos coprimos. Ponemos $n$ cartas en círculo,
una en cada una de las posiciones $0,1,\dots,n-1$, numeradas en sentido
horario. Un \emph{movimiento} es una permutación de las posiciones. Para cada
$k$ con $0\le k<n$, sea $a_k$ la cantidad de movimientos que, aplicados $m$
veces seguidas, corren cada carta $k$ lugares en sentido horario.

Sea $\omega$ un número complejo tal que $\omega^n=1$.

\begin{enumerate}
  \item\itshape Sea $T$ una permutación de un conjunto finito, con $T^n=\id$.
        Para cada $k$ con $0\le k<n$, sea $b_k$ la cantidad de elementos que
        $T^k$ deja fijos (acá $T^0=\id$). Probar que
        $\sum_{k=0}^{n-1}b_k\omega^k$ es un múltiplo no negativo de $n$.
  \item\itshape Probar que $\sum_{k=0}^{n-1}a_k\omega^k$ también es un
        múltiplo no negativo de $n$.
\end{enumerate}

\section*{Solución}

Hacemos primero (1), que sirve como hint. 
Miremos un ciclo de $T$ de longitud $d$. Aplicar $T$ es avanzar un lugar en el
ciclo, así que volvemos al punto de partida después de $k$ pasos justo cuando
$d\mid k$. Como $T^n=\id$, en particular $d\mid n$. Entonces el ciclo aporta
$d$ elementos a $b_k$ si $d\mid k$, y ninguno si no. Su aporte a la suma es
\[
  d\sum_{j=0}^{n/d-1}(\omega^d)^j
  =
  \begin{cases}
    n & \text{si } \omega^d=1,\\
    0 & \text{si } \omega^d\ne 1.
  \end{cases}
\]
Si $\omega^d=1$, hay $n/d$ sumandos y todos valen $1$. Si no, es una suma
geométrica que se anula, porque $(\omega^d)^{n/d}=\omega^n=1$.

Sumamos los aportes de todos los ciclos y listo: la suma es $n$ por la
cantidad de ciclos de $T$ cuya longitud $d$ cumple $\omega^d=1$.

Ahora resolvemos (2). 
Sea $R$ la rotación que corre cada carta un lugar en sentido horario. Fijemos
$k$. Un movimiento $\pi$ cuenta para $a_k$ si y solo si $\pi^m=R^k$.

Primero buscamos una rotación que, aplicada $m$ veces, haga lo mismo que
$R^k$. Como $m$ y $n$ son coprimos, hay un entero positivo $u$ tal que
$mu\equiv 1\pmod n$. Tomamos $B=(R^k)^u$. Como $R^n=\id$, vale que
$B^m=R^k$.

Si $\pi^m=R^k$, por supuesto $\pi$ conmuta con $R^k$, que es una potencia de
$\pi$. Y como $B$ es una potencia de $R^k$, $\pi$ también conmuta con $B$. Sea
$\tau=B^{-1}\pi$. Entonces $\tau$ conmuta con $R^k$ y
$\tau^m=B^{-m}\pi^m=\id$. Al revés, si $\tau^m=\id$ y $\tau$ conmuta con
$R^k$, también conmuta con $B$, y entonces $(B\tau)^m=B^m\tau^m=R^k$. Las dos
construcciones son una la inversa de la otra, porque $\pi=B\tau$. Por eso,
$a_k$ es la cantidad de permutaciones $\tau$ con $\tau^m=\id$ que conmutan con
$R^k$.

Falta elegir el conjunto y la permutación $T$. El conjunto es el de todas las
permutaciones $\tau$ con $\tau^m=\id$. La permutación $T$ renumera las
posiciones sumándoles uno, módulo $n$: si $\tau$ lleva $i$ a $j$, entonces
$T(\tau)$ lleva $i+1$ a $j+1$. En fórmulas, $T(\tau)=R\tau R^{-1}$.
Renumerar no cambia la condición $\tau^m=\id$, y después de $n$ pasos
volvemos a la numeración original. Así que $T$ es una permutación del conjunto
y $T^n=\id$. Además,
\[
  T^k(\tau)=\tau
  \quad\Longleftrightarrow\quad R^k\tau R^{-k}=\tau
  \quad\Longleftrightarrow\quad R^k\tau=\tau R^k.
\]
O sea, con este $T$ vale que $b_k=a_k$ para todo $k$. Alcanza con aplicar la
primera parte.

\section*{Observación}

Para que el problema quede más difícil, se puede
pedir directamente (2).

\end{document}
